\documentclass[../../main.tex]{subfiles}
\begin{document}


\begin{problem}
字符 a ~ h 的出现频率恰好是前 8 个斐波那契数。
\begin{itemize}
    \item 求他们的 Huffman 编码。
    \item 将结果推广到 n 个字符的频率恰好是前 n 个斐波那契数的情况。
\end{itemize}
\end{problem}
\begin{solution}

前 8 个斐波那契数为\verb|1,1,2,3,5,8,13,21|

Huffman哈夫曼编码树为:

\verb|                            [54]             |

\verb|                          0/   1\            |

\verb|                        [33]    [h:21]       |
  
\verb|                      0/   1\                |
      
\verb|                   [20]     [g:13]           |
        
\verb|                 0/   1\                     |
        
\verb|              [12]     [f:8]                 |
        
\verb|            0/   1\                          |
        
\verb|          [7]     [e:5]                      |
        
\verb|        0/  1\                               |
        
\verb|      [4]    [d:3]                           |
        
\verb|    0/  1\                                   |
        
\verb|   [2]   [c:2]                               |
        
\verb| 0/  1\                                      |
        
\verb|[a:1] [b:1]                                  |

所以:

a的编码为:\verb|0000000|

b的编码为:\verb|0000001|

c的编码为:\verb|000001|

d的编码为:\verb|00001|

e的编码为:\verb|0001|

f的编码为:\verb|001|

g的编码为:\verb|01|

h的编码为:\verb|1|

推广到n个字符:
\[
\begin{aligned}
&\text{第 1 个字符:} \quad \underbrace{0 \, 0 \, \cdots \, 0}_{n-1 \text{ 个 0}}, \\
&\text{第 2 个字符:} \quad \underbrace{0 \, 0 \, \cdots \, 0}_{n-2 \text{ 个 0}} \, 1, \\
&\text{第 3 个字符:} \quad \underbrace{0 \, 0 \, \cdots \, 0}_{n-3 \text{ 个 0}} \, 1, \\
&\quad \vdots \\
&\text{第 } (n-1) \text{ 个字符:}  0 \, 1, \\
&\text{第 n 个字符:}  1
\end{aligned}
\]
\end{solution}
\end{document}